\scanpage{24}
\mspage{12}
Inversement, donnons-nous \inserted{un groupe $T$ abélien} un $T$-groupoïde \deleted{\unsure{connexe}} $\Lambda$, une opération de $\mathcal E$ sur $(T,\Lambda)$, et un hom de groupoïdes
\diagram{\begin{tikzcd}[ampersand replacement=\&]
\Lambda \arrow[r,"\text{foncteur \uuline{fidèle}}"] \& \begin{gathered}\text{groupoïde à un seul objet}\\\text{des rev. universels de}\\\Pi\text{ égaux à }\widetilde\Pi\\\text{dont les autom sont }\pi\end{gathered}
\end{tikzcd}}
compatible avec action de $\mathcal E$.
Soit $\tau=\operatorname{Ob}\Lambda$ [donc $\tau$ est un $\mathcal E$-ens], et pour $\forall\,t,t'\in\tau$, on a un hom \uline{injectif}
\diagram{\begin{tikzcd}[ampersand replacement=\&]
\underbrace{\operatorname{Isom}_\Lambda(t,t')}_{\pi_{t,t'}} \arrow[r,hook] \& \pi
\end{tikzcd}}
compatible avec composition des isom.

On a exprimé la donnée d'un foncteur fidèle d'un groupoïde $\Lambda$ dans $\langle\pi\rangle$ --- \deleted{Mais} [il faut se donner une famille de sous groupes $(\pi_t)_{t\in\tau}$ de $\pi$ et \unsure{d'ens vides ou} bitorseurs sous les $\pi_{t'},\pi_t\ldots$] On va maintenant introduire les op. de
\mspage{23}
